% !PRINTED 609
\weberpage{609}
Jede ganze Zahl $\eta$ des Körpers $\mathfrak{K}$ kann (Bd.\,I, § 151) in der
Form dargestellt werden:
\[
\eta = \frac{\varphi(\xi)}{f'(\xi)},
\]
worin $\varphi(\xi)$ ganze Koeffizienten in $\Omega$ hat. Man kann dies in die
Form setzen:
\weq{9}{\gamma\eta = \gamma_0 + \gamma_1\xi + \gamma_2\xi^2 + \cdots + \gamma_{n-1}\xi^{n-1},}
worin $\gamma, \gamma_0, \gamma_1, \ldots, \gamma_{n-1}$ ganze Zahlen in $\Omega$ sind, und so, daß die
erste, $\gamma$, ein Teiler der Diskriminante von $f$ ist. Die Primfaktoren von
$\gamma$ schließen wir durch den Exkludenten aus.

Es sei

$\mathfrak{P}$ ein Primideal ersten Grades in $\mathfrak{K}$,

$\mathfrak{p}_1$ das durch $\mathfrak{P}$ teilbare Primideal ersten Grades in $\Omega$,

$p$ die Norm von $\mathfrak{p}_1$ und von $\mathfrak{P}$;

dann ist für jede ganze Zahl $\alpha$ in $\Omega$:
\[
\alpha^p \equiv \alpha \ (\mathrm{mod}\ \mathfrak{p}_1)
\]
und
\weq{10}{\xi^p \equiv \xi \ (\mathrm{mod}\ \mathfrak{P}),}
also nach (9) für jede ganze Zahl $\eta$ in $\mathfrak{K}$:
\[
\eta^p \equiv \eta \ (\mathrm{mod}\ \mathfrak{P}).
\]

9. \gesperrt{Die notwendige und hinreichende Bedingung dafür, daß $\mathfrak{K}(\xi)$ Klassenkörper sei, besteht also darin, daß für jedes in einem $\mathfrak{p}_1$ aufgehende Primideal $\mathfrak{P}$ und nur für dieses die Kongruenz (10) erfüllt sei.}

Es seien nun
\[
\mathfrak{K} = \Omega(\xi), \quad \mathfrak{K}' = \Omega(\xi')
\]
zwei Klassenkörper derselben Gruppe $A$. Zunächst ergibt sich leicht, daß sie \gesperrt{von gleichem Grade} sein müssen. Denn für beide muß die Bedingung (5) erfüllt sein. Sind also $n$ und $n'$ die Grade, so muß $(s-1)^{n'-n}$ für $s=1$ endlich und von Null verschieden sein, was nur möglich ist, wenn $n=n'$ ist.

Durch Zusammensetzung von $\mathfrak{K}$ und $\mathfrak{K}'$ leiten wir einen dritten Körper ab:
\[
\mathfrak{K}'' = \Omega(\xi, \xi'),
\]
von dem wir zeigen, daß er auch ein Klassenkörper ist.
